% ============================================================
% OP_anomaly_proposal_v2.tex
% Revision per GPT review of v1. NO preprint edits by this file.
% Changes vs v1: AN1 notation + explicit Witten phrasing;
% AN2 narrowed (rep-data primary, Fujikawa secondary);
% AN3 branch structure spelled out, exclusion rewording,
% B-L convention table, compactness caveat, separate placement;
% references page-verified (APS).
% ============================================================
\section*{Proposal v2: anomaly-cancellation package (OP-GUT6a sharpening)}

\subsection*{0. Scope (unchanged from v1)}
Sharpen \S anomaly\_architecture / \S anomaly\_constraints by
(i) displaying the arithmetic, (ii) upgrading B4 to a lemma,
(iii) classifying the hypercharge solution set with honest caveats.
No derivation claims; OP-GUT6a/6b/6c and OP-GUT3b remain Open.

\subsection*{1. Decomposition (approved structure)}
AN1 = Level~A arithmetic verification of inherited assignments
(verification, not derivation).
AN2 = candidate Level~A representation-data lemma at the local
anomaly-coefficient level.
AN3 = candidate Level~A algebraic classification under inherited
one-generation content.
AN4 = unchanged open remainder (OP-GUT6c, 't~Hooft, OP-GUT1/2/3b/5).

\subsection*{2. AN1: explicit verification (revised text core)}
With the assignments of eq.~hypercharges:
\begin{align*}
[SU(3)]^2 U(1)_Y&:\ 2\cdot\tfrac16-\tfrac23+\tfrac13=0,\\
[SU(2)]^2 U(1)_Y&:\ 3\cdot\tfrac16-\tfrac12=0,\\
U(1)_Y\text{--}\mathrm{grav}^2&:\
  6\cdot\tfrac16-3\cdot\tfrac23+3\cdot\tfrac13-2\cdot\tfrac12+1
  \;(+0)=1-2+1-1+1=0,\\
U(1)_Y^3&:\ 6(\tfrac16)^3+3(-\tfrac23)^3+3(\tfrac13)^3
  +2(-\tfrac12)^3+1^3\;(+0)
  =\tfrac1{36}-\tfrac{32}{36}+\tfrac{4}{36}-\tfrac{9}{36}
  +\tfrac{36}{36}=0.
\end{align*}
Witten $SU(2)$ anomaly: per generation there are $3$ colour copies
of $Q_L$ plus one $L_L$, hence $4N_{\rm gen}$ left-handed $SU(2)$
doublets, which is even.
Status: Level~A arithmetic verification of inherited assignments,
not a derivation; B3 scope labels (which conditions are
\emph{meaningful} in current ECT) are unchanged.

\subsection*{3. AN2: lemma (narrowed statement)}
\emph{Lemma (proposed).}
In the reconstructible ordered-branch chiral gauge EFT, the
deformation $\mathcal L_5=\mu_5\bar\Psi\gamma^\mu n_\mu\Psi$ does
not change the local gauge-anomaly coefficients computed from the
chiral representation data, and does not change the Witten $SU(2)$
doublet-counting class.
\emph{Proof (representation-data, primary).}
Anomaly coefficients are determined by the chiral fermion
representations and charges.
Since $\mathcal L_5$ changes neither the gauge representation nor
the chirality assignments --- it is the coupling of the background
vector $B_\mu=\mu_5 n_\mu$ to the gauge-singlet vector current
$\bar\Psi\gamma^\mu\Psi$ --- the anomaly coefficients are unchanged.
The $SU(2)$ doublet count is likewise unchanged.
\emph{Secondary check (Fujikawa, intuition only).}
Equivalently, in a Fujikawa description this vector background
shifts the Dirac operator without changing the chiral
representation trace that defines the anomaly coefficient.
\emph{Scope.}
Local gauge-anomaly coefficients and the Witten count only;
flavour-diagonal $\mathcal L_5$ as defined; possible mixed
gravitational/global or 't~Hooft-type bookkeeping involving the
background direction is outside the lemma; ECT application inherits
the B1 OS-closure conditionality.

\subsection*{4. AN3: solution structure (revised; own subsubsection)}
Proposed placement: separate block
\emph{``Solution structure of the one-generation anomaly
equations''} between Layer~C and the closing paragraphs of
\S anomaly\_architecture (not buried inside a Layer-C paragraph).

\emph{Proposition (candidate Level~A algebraic classification,
conditional on the inherited one-generation chiral content
$\{Q,u^c,d^c,L,e^c\}$, all left-handed Weyl, no $\nu^c$).}
The linear conditions give $Y_L=-3Y_Q$, $Y_{e^c}=6Y_Q$,
$Y_{u^c}+Y_{d^c}=-2Y_Q$; with $Y_{u^c}=-Y_Q+\delta$,
$Y_{d^c}=-Y_Q-\delta$ the cubic condition reduces to
\[
18\,Y_Q\,\bigl(9Y_Q^2-\delta^2\bigr)=0 .
\]
Apart from the trivial all-zero charge assignment, the solution set,
up to overall normalisation, consists of two nontrivial assignments:
(a) the \emph{non-degenerate branch} $Y_Q\neq0$,
$\delta=\pm3Y_Q$, where the two sign choices exchange the labels
$u^c\!\leftrightarrow\! d^c$ and reproduce the Standard-Model
assignment;
(b) the \emph{degenerate branch} $Y_Q=0$
($Y_L=Y_{e^c}=0$, $Y_{u^c}=-Y_{d^c}$), in which hypercharge acts
only on $u^c,d^c$ and the lepton sector is hypercharge-neutral.
This two-assignment structure matches the published classification
(Minahan--Ramond--Warner; Geng--Marshak), where branch (b) corresponds to the non-SM assignment sometimes
called ``bizarre'' in that literature.
\emph{Exclusion of the degenerate branch.}
Branch (b) is excluded by the inherited SM charge-assignment
requirement for the charged-lepton sector --- equivalently by the
inherited identification of $e^c$ as the charged-lepton singlet
with nonzero hypercharge.
Anomaly cancellation alone does not exclude it; this exclusion is
recorded explicitly as inherited input.
\emph{Corollary 1 (Part-IV scoping).}
The sentence ``these constraints fix hypercharges up to
normalisation'' in \S anomaly\_constraints is correct for the
listed content modulo the $u^c\!\leftrightarrow\! d^c$ relabelling
and the excluded degenerate branch; add this scoping with a
pointer.
\emph{Corollary 2 ($\nu^c$ and the $B{-}L$ flat direction).}
Starting from the non-degenerate SM-like branch, if a singlet $\nu^c$
with \emph{free} hypercharge is admitted, the shifted assignment
$Y\to Y+t\,(B{-}L)$ is anomaly-free for all $t$.
Here $B{-}L$ is understood in the left-handed conjugate convention:
$Q_L:+\tfrac13$,\quad $u^c,d^c:-\tfrac13$,\quad $L_L:-1$,\quad
$e^c,\nu^c:+1$.
Hence the anomaly system alone no longer fixes hypercharge.
If $\nu^c$ is admitted but constrained from the outset to be a
hypercharge-neutral singlet, $Y_{\nu^c}=0$, this particular
$B{-}L$ flat direction is removed; the existing Part-IV sentence ``a singlet can
be added without disturbing anomaly balance'' implicitly makes this
choice --- propose stating it explicitly.
\emph{Corollary 3 (compactness).}
The ECT compact-$U(1)$ lattice (Theorem~charge\_lattice;
$2T+Y\in2\mathbb Z$) enforces rationality/lattice structure but
does \emph{not} remove the $B{-}L$ flat direction (which is
rational-valued), unless additional normalisation or
charge-identification input is imposed.
The observed assignment therefore remains OP-GUT3b-level input.

\subsection*{5. Implementation map (after approval only)}
\begin{enumerate}
  \item Layer~B: new paragraph ``B2$'$. Explicit verification for
    the inherited assignments'' (AN1) after B2.
  \item B4: insert the AN2 lemma (narrowed statement, rep-data
    proof, Fujikawa as secondary check); keep surrounding prose.
  \item New \emph{subsubsection}
    ``Solution structure of the one-generation anomaly equations''
    (AN3 + corollaries) between Layer~C and the closing paragraphs.
  \item Status table: rows for AN1 (A, verification of inherited
    assignments), AN2 (A candidate, representation-data lemma),
    AN3 (A candidate, algebraic classification under inherited
    content); adjust the $\mathcal L_5$ row comment.
  \item \S anomaly\_constraints: one-sentence scoping of the
    uniqueness claim + explicit $Y_{\nu^c}=0$ remark + pointer.
  \item OP-GUT6a/6b rows: comments gain ``explicit consistency
    verification recorded for the inherited content''; statuses
    remain Open as derivations.
  \item Consistency sweep at implementation: grand tables, Part-I
    derivation graph (anomaly-node check), abbreviations (ABJ
    check), companion deferred (Q11 pattern).
\end{enumerate}

\subsection*{6. References (page-verified, APS abstract pages)}
\begin{itemize}
  \item C.~Q.~Geng and R.~E.~Marshak, ``Uniqueness of quark and
    lepton representations in the standard model from the anomalies
    viewpoint,'' Phys.\ Rev.\ D \textbf{39}, 693 (1989).
    [DOI 10.1103/PhysRevD.39.693 --- verified]
  \item J.~A.~Minahan, P.~Ramond and R.~C.~Warner, ``Comment on
    anomaly cancellation in the standard model,'' Phys.\ Rev.\ D
    \textbf{41}, 715 (1990).
    [DOI 10.1103/PhysRevD.41.715 --- verified; abstract states the
    two-assignment result used in AN3]
  \item Optional third (decide at review): Geng--Marshak Reply,
    Phys.\ Rev.\ D \textbf{41}, 717 (1990) [verified], discussing
    the rejection of the ``bizarre'' assignment.
\end{itemize}
Proposal: include the first two only.

\subsection*{7. Status discipline (unchanged)}
Forbidden: ``ECT derives anomaly cancellation'';
``anomalies (+ compactness) fix the observed hypercharges'';
``ECT predicts SM matter content''.
Safe: ``verified for inherited assignments'';
``classification under inherited one-generation content'';
``consistency test passed''.
